\begin{abstract}
배틀로얄 장르의 게임에서 최종 순위는 경기 종료 후 기록되는 단일 결과값이지만,
실제로는 각 팀이 시간에 따라 변화하는 공간, 자원, 교전 사건, 생존 상태 속에서
여러 페이즈를 통과하거나 탈락하는 과정의 결과다. 이 논문은 PUBG: Battlegrounds
스쿼드 매치의 공식 텔레메트리(telemetry)를 경기 후 집계 행으로 보지 않고,
시공간 데이터베이스 표현으로 다시 구성한다. 스쿼드 팀은 이동 객체, 안전구역과
블루존은 시간변화 영역, 피해 사건, 기절 사건, 처치 사건, 탈락 사건은 시공간
사건으로 둔다. 이를 페이즈별 위험집합(risk set)과 10초 간격 팀 상태 관측값으로
변환한다.

최종 순위 예측은 이 표현 위에서 다음 페이즈 탈락 위험도(hazard) 예측 문제로
바뀐다. 9개 맵/플레이 모드 조건의 2,500개 스쿼드 매치로 초기 비교 실험을 수행한
결과, 가장 안정적인 신호는 페이즈별로 갱신되는 팀 생존 여력이었다. 이는 전투 생존
게임에서 평균 체력, 최소 체력, 생존 인원, 차량 이용 여부 같은 직접 상태 변수가
강하게 작동한 결과로 해석해야 한다. 더 중요한 비교는 직접 생존 상태를 제외한
조건에서도 이동 이력이 의미를 갖는지이다. 이동 이력
조건은 팀 중심 좌표와 속도의 10초 이력만으로 다음 페이즈 탈락을 예측했으며, 페이즈
종료 시점의 단일 상태만 사용하는 기준모델(baseline)보다 낮은 손실을 보였다. 반대로
안전구역 맥락은 단독 조건과 이동 이력과의 단순 결합 조건에서 뚜렷한 개선을 만들지
못했다.

이 결과를 안전구역이 불필요하다는 뜻으로 읽기는 어렵다. 더 정확한 해석은 현재
구성에서 먼저 보존되어야 할 정보가 팀의 생존 상태와 이동 이력이며, 안전구역은
독립 특징으로 붙이는 방식보다 이동 이력과 생존 상태의 의미를 조건화하는 표현으로
다시 설계되어야 한다는 것이다.
\end{abstract}
